\chapter{Relationale Abbildungen und endliche Abschlüsse}
\label{chap:b37-guide-relational}

\section{Gruppenkern und schwache Konjugation}

\begin{definition}[Relationaler Morphismus]
\label{def:b37-guide-relational-morphism}
Ein \emph{relationaler Morphismus} $\varphi:A\relbar\joinrel\to B$
von einer Halbgruppe $A$ in eine Halbgruppe $B$ ordnet jedem $a\in A$
eine nichtleere Teilmenge $\varphi(a)\subseteq B$ zu, so daß
\[
 \varphi(a)\varphi(b)\subseteq\varphi(ab)\qquad(a,b\in A)
\]
gilt. Dabei bezeichnet das Produkt zweier Teilmengen die Menge aller
Produkte ihrer Elemente. Für $t\in B$ schreiben wir
\[
 \varphi^{-1}(t)=\{a\in A:t\in\varphi(a)\}.
\]
\end{definition}

\begin{definition}[Gruppenkern und Abschluß unter schwacher Konjugation]
\label{def:b37-guide-group-kernel}
Der \emph{Gruppenkern} von $S$ ist die Menge
\[
 K_{\mathsf G}(S)=\bigcap_{G,\varphi}\varphi^{-1}(1_G),
\]
wobei $G$ alle endlichen Gruppen und $\varphi$ alle relationalen
Morphismen $S\relbar\joinrel\to G$ durchläuft.
Das Symbol $K_{\mathsf G}(S)$ bezeichnet nicht das kleinste Ideal $K(S)$.

Eine Teilmenge $C\subseteq S$ heißt \emph{unter schwacher Konjugation
abgeschlossen}, wenn
\[
 u\in C,\quad x,y\in S,\quad xyx=x
 \quad\Longrightarrow\quad xuy\in C\text{ und }yux\in C
\]
gilt. Mit $C_{\mathsf G}(S)$ bezeichnen wir die kleinste
produktabgeschlossene Teilmenge von $S$, die $E(S)$ enthält und unter
schwacher Konjugation abgeschlossen ist.
Für $S=\varnothing$ sind beide Mengen leer.
\end{definition}

\begin{theorem}[Typ-II-Satz von Ash: Vollständigkeitsrichtung]
\label{thm:b37-guide-type-two}
Für jede endliche Halbgruppe $S$ gilt
\[
 K_{\mathsf G}(S)\subseteq C_{\mathsf G}(S).
\]
\cite{Ash1991,RhodesSteinberg2009}
\end{theorem}

\section{Gruppenpunktartige Teilmengen}

\begin{definition}[Gruppenpunktartige Teilmenge]
Eine nichtleere Teilmenge $X\subseteq S$ heißt \emph{gruppenpunktartig},
wenn für jede endliche Gruppe $G$ und jeden relationalen Morphismus
$\varphi:S\relbar\joinrel\to G$ ein $g\in G$ existiert mit
\[
 X\subseteq\varphi^{-1}(g).
\]
Die Familie dieser Teilmengen wird mit
$\operatorname{PL}_{\mathsf G}(S)$ bezeichnet.
\end{definition}

\begin{definition}[Schwache Inversenmengen]
Wir setzen $M=S^1$ und $D=C_{\mathsf G}(M)$.
Für $m\in M$ seien
\[
 m^{(+)}=\{m\},\qquad
 m^{(-)}=\{x\in M:xmx=x\}.
\]
Die zweite Menge ist die Menge der \emph{schwachen Inversen} von $m$.
\end{definition}

\begin{theorem}[Gruppenpunktartige Teilmengen nach Ash: Vollständigkeitsrichtung]
\label{thm:b37-guide-group-pointlikes}
Zu jeder Teilmenge $X\in\operatorname{PL}_{\mathsf G}(S)$ gibt es
eine ganze Zahl $r\geq1$, Elemente $m_1,\ldots,m_r\in M$ und Zeichen
$\varepsilon_1,\ldots,\varepsilon_r\in\{+,-\}$ mit
\[
 X\subseteq Dm_1^{(\varepsilon_1)}D\cdots
                 m_r^{(\varepsilon_r)}D.
\]
\cite[Satz~1.2]{HenckellMargolisPinRhodes1991}
\end{theorem}

\clearpage
\section{Aperiodisch punktartige Teilmengen}

\begin{definition}[Potenzhalbgruppe und zyklischer Mengenabschluß]
\label{def:b37-guide-power-subsets}
Es sei
\[
 \mathcal P_+(S)=\{X\subseteq S:X\ne\varnothing\}
\]
die Menge der nichtleeren Teilmengen von $S$, versehen mit dem
\emph{Mengenprodukt}
\[
 XY=\{xy:x\in X,\ y\in Y\}.
\]
Für $X\in\mathcal P_+(S)$ werden $X^k$ und $X^\omega$ in dieser
endlichen Halbgruppe gebildet. Insbesondere ist $X^\omega$ die
idempotente Potenz der Teilmenge $X$ bezüglich des Mengenprodukts.
Wir setzen
\[
 X^{\omega+*}=X^\omega\ \cup\!\bigcup_{k\geq1}X^\omega X^k.
\]
\end{definition}

\begin{definition}[Aperiodisch punktartige Teilmenge]
\label{def:b37-guide-aperiodic-pointlike}
Eine nichtleere Teilmenge $X\subseteq S$ heißt
\emph{aperiodisch punktartig}, wenn für jede nichtleere endliche
aperiodische Halbgruppe $T$ und jeden relationalen Morphismus
$\varphi:S\relbar\joinrel\to T$ ein $t\in T$ existiert mit
\[
 X\subseteq\varphi^{-1}(t).
\]
Die Familie aller aperiodisch punktartigen Teilmengen von $S$
wird mit $\operatorname{PL}_{\mathsf A}(S)$ bezeichnet.
\end{definition}

\begin{definition}[Abschlußfamilie für aperiodische Punktartigkeit]
\label{def:b37-guide-aperiodic-closure}
Mit $\mathcal C_{\mathsf A}(S)$ bezeichnen wir die kleinste Familie
$\mathcal C\subseteq\mathcal P_+(S)$, die alle Einermengen
$\{s\}$ mit $s\in S$ enthält und unter Mengenprodukten, nichtleeren
Teilmengen und der Abbildung $X\mapsto X^{\omega+*}$ abgeschlossen ist.
\end{definition}

\begin{theorem}[Abschlußsatz von Henckell: Vollständigkeitsrichtung]
\label{thm:b37-guide-aperiodic-pointlikes}
Für jede endliche Halbgruppe $S$ gilt
\[
 \operatorname{PL}_{\mathsf A}(S)\subseteq\mathcal C_{\mathsf A}(S).
\]
\cite{Henckell1988,VanGoolSteinberg2019}
\end{theorem}

\section{Idempotente aperiodische Punktartigkeit}

\begin{definition}[Idempotent aperiodisch punktartige Teilmenge]
Eine nichtleere Teilmenge $X\subseteq S$ heißt
\emph{idempotent aperiodisch punktartig}, wenn für jede nichtleere
endliche aperiodische Halbgruppe $T$ und jeden relationalen Morphismus
$\varphi:S\relbar\joinrel\to T$ ein $e\in E(T)$ existiert mit
\[
 X\subseteq\varphi^{-1}(e).
\]
\end{definition}

\begin{theorem}[Idempotente Punktartigkeit nach Henckell: Existenzrichtung]
\label{thm:b37-guide-idempotent-pointlikes}
Zu jeder idempotent aperiodisch punktartigen Teilmenge $X\subseteq S$
gibt es ein $Y\in\operatorname{PL}_{\mathsf A}(S)$ mit
\[
 X\subseteq Y,\qquad Y^2=Y.
\]
Die Gleichung $Y^2=Y$ bezieht sich auf das Mengenprodukt.
\cite{Henckell2004,HenckellRhodesSteinberg2010StablePairs}
\end{theorem}

\section{Aperiodisch stabile Paare}

\begin{definition}[Aperiodisch stabiles Paar]
Ein Paar $(X,U)$ aus einer nichtleeren Teilmenge $X\subseteq S$
und einer Unterhalbgruppe $U\subseteq S$ heißt \emph{aperiodisch stabil},
wenn für jeden relationalen Morphismus
$\varphi:S\relbar\joinrel\to T$ in eine nichtleere endliche
aperiodische Halbgruppe ein $t\in T$ existiert mit
\[
 X\subseteq\varphi^{-1}(t),\qquad
 \varphi(u)\cap\operatorname{Stab}_{R}^{T}(t)\ne\varnothing
 \quad(u\in U).
\]
\end{definition}

\begin{definition}[Interne $\mathcal L$-Kette von Teilmengen]
Eine nichtleere, unter Mengenprodukten abgeschlossene Familie
$\mathcal W\subseteq\mathcal P_+(S)$ heißt
\emph{interne $\mathcal L$-Kette}, wenn für alle $A,B\in\mathcal W$
mindestens eine der folgenden Bedingungen gilt:
\[
 A=B,\qquad
 A=CB\text{ für ein }C\in\mathcal W,\qquad
 B=CA\text{ für ein }C\in\mathcal W.
\]
\end{definition}

\begin{theorem}[Stabile Paare nach Henckell: Existenzrichtung]
\label{thm:b37-guide-stable-pairs}
Zu jedem aperiodisch stabilen Paar $(X,U)$ einer endlichen Halbgruppe
$S$ gibt es ein $Y\in\operatorname{PL}_{\mathsf A}(S)$ und eine
interne $\mathcal L$-Kette
$\mathcal W\subseteq\operatorname{PL}_{\mathsf A}(S)$ mit
\[
 X\subseteq Y,\qquad
 U\subseteq\bigcup_{A\in\mathcal W}A,\qquad
 YA=Y\quad(A\in\mathcal W).
\]
\cite[Satz~6.4]{HenckellRhodesSteinberg2010StablePairs}
\end{theorem}
